\documentclass[11pt,a4paper]{article}
\usepackage[utf8]{inputenc}
\usepackage[T1]{fontenc}
\usepackage{lmodern}
\usepackage[margin=2.2cm]{geometry}
\usepackage{amsmath,amssymb,amsthm,mathtools}
\usepackage{enumitem}
\usepackage{xcolor}
\usepackage[most]{tcolorbox}
\usepackage{tikz}
\usetikzlibrary{arrows.meta,calc}
\usepackage{fancyhdr}
\usepackage[hidelinks]{hyperref}
\usepackage{multicol}

\definecolor{azul}{RGB}{40,70,150}
\definecolor{verde}{RGB}{30,120,70}
\definecolor{rojo}{RGB}{170,40,40}
\definecolor{gris}{RGB}{90,90,90}

\setlength{\parindent}{0pt}
\setlength{\parskip}{5pt}

\newtcolorbox{resumen}[1][]{colback=azul!5,colframe=azul,fonttitle=\bfseries,title={#1},breakable,arc=2pt,boxrule=0.8pt}
\newtcolorbox{receta}[1][Receta]{colback=verde!6,colframe=verde,fonttitle=\bfseries,title={#1},breakable,arc=2pt,boxrule=0.8pt}
\newtcolorbox{ojo}[1][Errores típicos]{colback=rojo!5,colframe=rojo,fonttitle=\bfseries,title={#1},breakable,arc=2pt,boxrule=0.8pt}
\newtcolorbox{ejemplo}[1][]{colback=white,colframe=gris,fonttitle=\bfseries,title={#1},breakable,arc=2pt,boxrule=0.6pt}

\newcounter{ejer}
\newcommand{\ejercicio}[2][]{\refstepcounter{ejer}\par\medskip\noindent\textbf{Ejercicio \theejer.}\ \textit{\small #2}\ifx&#1&\else\hfill{\small\textcolor{gris}{[#1]}}\fi\par}
\newcommand{\sol}[1]{\par\medskip\noindent\textbf{Ejercicio #1.}\ }

\newcommand{\R}{\mathbb{R}}
\newcommand{\dA}{\,dA}
\newcommand{\dV}{\,dV}

\pagestyle{fancy}
\fancyhf{}
\renewcommand{\headrulewidth}{0.4pt}
\fancyfoot[C]{\small\thepage}
\renewcommand{\figurename}{Figura}
\renewcommand{\tablename}{Tabla}
\fancyhead[L]{\small Cálculo para Ciencia de Datos -- IMT2220}
\fancyhead[R]{\small Guía 5: Análisis de Fourier}

\begin{document}

\begin{center}
{\LARGE\bfseries Guía 5: Análisis de Fourier}\\[6pt]
{\large Cálculo para Ciencia de Datos (IMT2220) -- Preparación Examen 2026-2}\\[3pt]
{\small Material: Clases 23, 24 y 25 \quad$\cdot$\quad Exámenes 2022-1, 2023-1 y 2025-2}
\end{center}

\section{Resumen}

\begin{resumen}[Señales periódicas]
$s$ tiene período $T$ si $s(t+T)=s(t)$. Su frecuencia es $f=\frac1T$ (Hz) y su frecuencia angular es $\omega=2\pi f=\frac{2\pi}T$.
\begin{itemize}[leftmargin=*,itemsep=1pt]
\item Una suma de sinusoides con frecuencias enteras $f_1,\dots,f_n$ tiene período $1/\mathrm{mcd}(f_1,\dots,f_n)$. Es periódica si y solo si los cocientes $f_i/f_j$ son racionales.
\item \textbf{Amplitud y fase:} $A\sin(\omega t+\varphi)=a\cos\omega t+b\sin\omega t$, con $a=A\sin\varphi$, $b=A\cos\varphi$ y $A=\sqrt{a^2+b^2}$. Cada frecuencia es un subespacio de dimensión 2: $\mathrm{span}\{\cos\omega t,\sin\omega t\}$.
\end{itemize}
\end{resumen}

\begin{resumen}[Producto interno de funciones y ortogonalidad]
\[
\langle f,g\rangle=\int_{-\pi}^{\pi}f(x)g(x)\,dx,\qquad \|f\|^2=\int_{-\pi}^\pi f^2\ \ (\text{energía}),\qquad \cos\theta=\frac{\langle f,g\rangle}{\|f\|\,\|g\|}.
\]
\textbf{Simetría:} si $f$ es impar y $g$ es par, $\langle f,g\rangle=0$.

\textbf{Teorema de ortogonalidad.} $\{1,\cos x,\sin x,\cos2x,\sin2x,\dots\}$ es ortogonal en $[-\pi,\pi]$, con
\[
\|1\|^2=2\pi,\qquad\|\cos kx\|^2=\|\sin kx\|^2=\pi .
\]
Se demuestra con las identidades producto-suma, por ejemplo $\sin\alpha\sin\beta=\frac12[\cos(\alpha-\beta)-\cos(\alpha+\beta)]$, y con $\int_{-\pi}^\pi\cos nx\,dx=0$ para $n\neq0$.
\end{resumen}

\begin{resumen}[Serie de Fourier (período $2\pi$)]
\[
f(x)\sim a_0+\sum_{k\ge1}\big(a_k\cos kx+b_k\sin kx\big)
\]
\[
\boxed{a_0=\frac1{2\pi}\int_{-\pi}^\pi f,\qquad a_k=\frac1\pi\int_{-\pi}^\pi f(x)\cos kx\,dx,\qquad b_k=\frac1\pi\int_{-\pi}^\pi f(x)\sin kx\,dx}
\]
$a_0$ es el promedio de $f$. Los coeficientes son las coordenadas de la \textbf{proyección ortogonal} $c=\frac{\langle f,\phi\rangle}{\langle\phi,\phi\rangle}$.
La suma parcial $S_nf$ es la \textbf{mejor aproximación} (mínimos cuadrados) de $f$ en $V_n=\mathrm{span}\{1,\cos x,\dots,\sin nx\}$, y el error $f-S_nf$ es ortogonal a $V_n$.

\textbf{Atajo:} si $f$ es \textbf{par}, $b_k=0$ y $a_k=\frac2\pi\int_0^\pi f\cos kx$ (solo cosenos). Si $f$ es \textbf{impar}, $a_k=0$ y $b_k=\frac2\pi\int_0^\pi f\sin kx$ (solo senos).

\textbf{Período $T$:} se reemplaza $kx$ por $\frac{2\pi kt}{T}$, con $a_0=\frac1T\int_0^Tf$, $a_k=\frac2T\int_0^Tf\cos\frac{2\pi kt}T$ y $b_k=\frac2T\int_0^Tf\sin\frac{2\pi kt}T$. El término $k$ tiene frecuencia $k/T$ Hz.

\textbf{Ojo con la convención:} en la Clase 24, $a_0$ es el promedio. Algunas pautas antiguas escriben $\frac{a_0}2$ con $a_0=\frac1\pi\int f$: es la misma serie.
\end{resumen}

\begin{resumen}[Convergencia (Dirichlet) y espectro]
Si $f$ es $2\pi$-periódica y $C^1$ a trozos:
\[
\lim_{n\to\infty}S_nf(x)=\frac{f(x^+)+f(x^-)}2 .
\]
Donde $f$ es continua converge a $f(x)$, y en un salto converge al \textbf{promedio de los límites laterales}. Cerca de los saltos queda una oscilación que no baja (fenómeno de Gibbs).

\textbf{Espectro de amplitudes:} $A_k=\sqrt{a_k^2+b_k^2}$ contra la frecuencia. Si $f$ tiene saltos, $A_k\sim1/k$; si tiene esquinas, $A_k\sim1/k^2$; si es suave, decae muy rápido. Una señal más suave necesita menos términos.
\end{resumen}

\begin{resumen}[Teorema de Parseval (Pitágoras en dimensión infinita)]
\[
\frac1{2\pi}\int_{-\pi}^{\pi}f^2\,dx=a_0^2+\frac12\sum_{k\ge1}\big(a_k^2+b_k^2\big)=a_0^2+\sum_{k\ge1}\frac{A_k^2}2 .
\]
La potencia de la señal es la suma de las potencias de sus armónicos (cada armónico aporta $\frac12A_k^2$, sin importar la fase).
\begin{itemize}[leftmargin=*,itemsep=1pt]
\item \textbf{Error de truncar:} $\frac1{2\pi}\int(f-S_nf)^2=\frac12\sum_{k>n}(a_k^2+b_k^2)$, la energía descartada.
\item \textbf{Con $m$ armónicos}, lo óptimo es guardar los $m$ de \emph{mayor energía} (no necesariamente los primeros).
\item Sirve para sumar series: $\sum\frac1{k^2}=\frac{\pi^2}6$, $\sum_{k\text{ impar}}\frac1{k^2}=\frac{\pi^2}8$, $\sum\frac1{k^4}=\frac{\pi^4}{90}$.
\end{itemize}
\end{resumen}

\begin{resumen}[Muestreo y aliasing]
Muestras $s_n=s(n\Delta t)$, con frecuencia de muestreo $f_s=\frac1{\Delta t}$. Las frecuencias $f$, $f+mf_s$ y $mf_s-f$ dan las \textbf{mismas muestras}.
La \textbf{frecuencia aparente} está en $[0,f_s/2]$: se resta el múltiplo de $f_s$ más cercano y se toma valor absoluto. Si $f<f_s/2$ no hay aliasing.
\end{resumen}

\begin{resumen}[Tres series que conviene saber (período $2\pi$)]
\begin{tabular}{@{}lll@{}}
Onda cuadrada ($-1$ en $(-\pi,0)$, $1$ en $(0,\pi)$) & impar & $\displaystyle\frac4\pi\Big(\sin x+\frac{\sin3x}3+\frac{\sin5x}5+\cdots\Big)$\\[8pt]
Diente de sierra $f=x$ en $(-\pi,\pi)$ & impar & $\displaystyle2\sum_{k\ge1}\frac{(-1)^{k+1}}{k}\sin kx$\\[8pt]
$f=|x|$ en $[-\pi,\pi]$ & par & $\displaystyle\frac\pi2-\frac4\pi\Big(\cos x+\frac{\cos3x}9+\frac{\cos5x}{25}+\cdots\Big)$
\end{tabular}
\end{resumen}

\begin{ojo}
\begin{itemize}[leftmargin=*,itemsep=1pt]
\item Olvidar que $\cos k\pi=(-1)^k$ y $\sin k\pi=0$.
\item No aprovechar la paridad (se duplica el trabajo).
\item En funciones definidas por tramos, integrar en todo $[-\pi,\pi]$ en vez de solo donde $f\neq0$.
\item Decir que la serie converge a $f$ en un salto: converge al promedio.
\item Para $a_0$, usar $\frac1\pi$ en lugar de $\frac1{2\pi}$ (con la convención del curso).
\end{itemize}
\end{ojo}

\section{Ejemplos resueltos}

\begin{ejemplo}[Ejemplo 1 (Clase 24). Onda cuadrada]
$f=-1$ en $(-\pi,0)$ y $f=1$ en $(0,\pi)$. Es impar, así que $a_k=0$ y
\[
b_k=\frac2\pi\int_0^\pi\sin kx\,dx=\frac2{\pi k}\big(1-\cos k\pi\big)=\frac2{\pi k}\big(1-(-1)^k\big)=\begin{cases}\frac4{\pi k}&k\text{ impar}\\0&k\text{ par.}\end{cases}
\]
En $x=\frac\pi2$, $f$ es continua y vale 1. Dirichlet da $1=\frac4\pi(1-\frac13+\frac15-\cdots)$, es decir $\frac\pi4=1-\frac13+\frac15-\cdots$ (Leibniz). En $x=0$ la serie vale $0=\frac{1+(-1)}2$.
\end{ejemplo}

\begin{ejemplo}[Ejemplo 2 (Clase 24 y Examen 2025-2). $f(x)=|x|$]
Es par: $b_k=0$, $a_0=\frac1\pi\int_0^\pi x\,dx=\frac\pi2$. Por partes,
\[
a_k=\frac2\pi\int_0^\pi x\cos kx\,dx=\frac2\pi\Big(\Big[\frac{x\sin kx}k\Big]_0^\pi-\int_0^\pi\frac{\sin kx}k\,dx\Big)=\frac2{\pi k^2}\big((-1)^k-1\big),
\]
es decir $a_k=-\frac4{\pi k^2}$ si $k$ es impar y $0$ si es par. Los coeficientes decaen como $1/k^2$ porque $|x|$ es continua (solo tiene esquinas).
\end{ejemplo}

\begin{ejemplo}[Ejemplo 3 (Clase 25). Parseval con la onda cuadrada]
$\frac1{2\pi}\int f^2=1$ y $b_k^2=\frac{16}{\pi^2k^2}$ ($k$ impar). Parseval: $1=\frac12\cdot\frac{16}{\pi^2}\sum_{k\text{ impar}}\frac1{k^2}$, así que $\sum_{k\text{ impar}}\frac1{k^2}=\frac{\pi^2}8$.
\end{ejemplo}

\section{Ejercicios}

\ejercicio[Clase 23]{Períodos, fase y ángulos}
\begin{enumerate}[label=(\alph*),itemsep=0pt]
\item Determine el período de $\sin(2\pi\cdot150t)+\sin(2\pi\cdot250t)$ ($t$ en segundos) y el de $\cos4t+\sin6t$.
\item Escriba $3\sin(t+\frac\pi3)$ en la forma $a\cos t+b\sin t$, y $\cos t-\sin t$ en la forma $A\sin(t+\varphi)$.
\item Calcule el ángulo entre $f(x)=x$ y $g(x)=1$, y entre $f(x)=|x|$ y $g(x)=1$, en $[-\pi,\pi]$.
\item Calcule $\langle\sin2x,\sin3x\rangle$ con una identidad producto-suma.
\end{enumerate}

\ejercicio[Clase 24]{Sin calcular integrales}
\begin{enumerate}[label=(\alph*),itemsep=0pt]
\item Encuentre la serie de Fourier de $\sin^3x$ (use $\sin3x=3\sin x-4\sin^3x$). ¿Para qué $n$ se cumple $S_nf=f$?
\item Use la ortogonalidad para calcular $\int_{-\pi}^\pi(2\sin x+\cos3x)^2dx$.
\item Proyecte $f(x)=x$ sobre $\sin x$: calcule $c=\langle x,\sin x\rangle/\langle\sin x,\sin x\rangle$.
\end{enumerate}

\ejercicio[Clase 24]{La serie de $x^2$}
Sea $f$ $2\pi$-periódica con $f(x)=x^2$ en $[-\pi,\pi]$.
\begin{enumerate}[label=(\alph*),itemsep=0pt]
\item Explique por qué $b_k=0$ y calcule $a_0$ y $a_k$.
\item Evalúe la serie en $x=\pi$ (justifique con Dirichlet) y deduzca $\sum_{k\ge1}\frac1{k^2}$.
\end{enumerate}

\ejercicio[Examen 2023-1]{Función por tramos}
Sea $g$ $2\pi$-periódica con $g(x)=0$ si $-\pi<x<0$ y $g(x)=x^2$ si $0\le x\le\pi$. Dibuje $g$ y encuentre su serie de Fourier. ¿A qué valor converge la serie en $x=\pi$?

\ejercicio[Examen 2022-1]{Función por tramos}
Sea $g$ $2\pi$-periódica con $g(x)=0$ si $-\pi<x<\frac\pi2$ y $g(x)=x$ si $\frac\pi2\le x\le\pi$. Dibuje $g$ y encuentre su serie de Fourier. ¿A qué valores converge la serie en $x=\frac\pi2$ y en $x=\pi$?

\ejercicio[Clase 24]{Señal digital}
Una señal vale 1 en $[0,\frac T2)$ y 0 en $[\frac T2,T)$, con $T=1$ ms. Escriba $f=\frac12(1+q)$, con $q$ una onda cuadrada, y obtenga su serie. Dé las frecuencias (en Hz) y amplitudes de sus tres primeros armónicos no nulos.

\ejercicio[Clase 25]{Parseval}
\begin{enumerate}[label=(\alph*),itemsep=0pt]
\item Para $g(x)=2\cos3x-4\sin3x$: ¿cuánto vale $A_3$? ¿Y la potencia media $\frac1{2\pi}\int g^2$?
\item Una temperatura es $f(t)=3+2\cos\frac{2\pi t}{24}+\sin\frac{4\pi t}{24}$ ($t$ en horas). Calcule $\frac1{24}\int_0^{24}f^2$ con Parseval. ¿Qué fracción de la potencia corresponde a las oscilaciones?
\item ¿Qué fracción de la energía de la onda cuadrada pierde $S_5f$?
\item Aplique Parseval a $|x|$ para obtener $\sum_{k\text{ impar}}\frac1{k^4}$ y deduzca $\sum_{k\ge1}\frac1{k^4}$.
\end{enumerate}

\ejercicio[Clase 25]{Guardar pocos armónicos}
Una señal tiene amplitudes $A_1=0.6$, $A_2=2$, $A_4=0.2$, $A_7=1.2$ y las demás nulas. Si solo se pueden guardar 2 armónicos, ¿cuáles conviene guardar? ¿Qué error relativo se comete? (Use $\sqrt{E_{\text{descartada}}/E_{\text{total}}}$.)

\ejercicio[Clase 25]{Aliasing}
\begin{enumerate}[label=(\alph*),itemsep=0pt]
\item Con $f_s=1000$ Hz, ¿hay aliasing para sinusoides de 450, 700, 1300 y 2000 Hz? Dé la frecuencia aparente en cada caso.
\item Un ECG se muestrea a $f_s=250$ Hz y hay interferencia muscular de 300 Hz. ¿En qué frecuencia aparece? ¿Por qué es molesto en Chile, donde la red eléctrica es de 50 Hz?
\item Una rueda con un rayo marcado gira a 26 vueltas/s y se filma a 24 cuadros/s. ¿Cómo se ve en pantalla?
\end{enumerate}

\newpage
\section{Soluciones}

\sol{1}
(a) $\mathrm{mcd}(150,250)=50$, así que el período es $\frac1{50}$ s $=20$ ms. Para $\cos4t+\sin6t$, los períodos son $\frac\pi2$ y $\frac\pi3$, y el período común es $\pi$.
(b) $3\sin(t+\frac\pi3)=\frac{3\sqrt3}2\cos t+\frac32\sin t$. Para $\cos t-\sin t$: $a=1$, $b=-1$, $A=\sqrt2$, $\varphi=\frac{3\pi}4$, es decir $\sqrt2\sin(t+\frac{3\pi}4)$.
(c) $x$ es impar y $1$ es par, así que el ángulo es de $90^\circ$. Para $|x|$: $\langle|x|,1\rangle=\pi^2$, $\|\,|x|\,\|=\sqrt{2\pi^3/3}$, $\|1\|=\sqrt{2\pi}$, luego $\cos\theta=\frac{\sqrt3}2$ y $\theta=30^\circ$.
(d) $\sin2x\sin3x=\frac12(\cos x-\cos5x)$, cuya integral en $[-\pi,\pi]$ es $0$.

\sol{2}
(a) $\sin^3x=\frac34\sin x-\frac14\sin3x$: $b_1=\frac34$, $b_3=-\frac14$ y el resto es $0$. $S_nf=f$ para $n\ge3$.
(b) Por Pitágoras: $4\|\sin x\|^2+\|\cos3x\|^2=4\pi+\pi=5\pi$.
(c) $\langle x,\sin x\rangle=2\pi$ (por partes) y $\langle\sin x,\sin x\rangle=\pi$, así que $c=2$ (coincide con $b_1$ del diente de sierra).

\sol{3}
(a) $f$ es par. $a_0=\frac1\pi\int_0^\pi x^2=\frac{\pi^2}3$ y $a_k=\frac2\pi\int_0^\pi x^2\cos kx\,dx=\frac{4(-1)^k}{k^2}$ (por partes dos veces).
(b) $f$ es continua en $\pi$ (porque $f(-\pi)=f(\pi)$): $\pi^2=\frac{\pi^2}3+\sum\frac{4(-1)^k(-1)^k}{k^2}=\frac{\pi^2}3+4\sum\frac1{k^2}$, así que $\sum\frac1{k^2}=\frac{\pi^2}6$.

\sol{4}
Solo se integra en $[0,\pi]$. $a_0=\frac1{2\pi}\int_0^\pi x^2dx=\frac{\pi^2}6$.
Por partes, $\int_0^\pi x^2\cos kx\,dx=\frac{2\pi(-1)^k}{k^2}$, así que $a_k=\frac{2(-1)^k}{k^2}$.
Además $\int_0^\pi x^2\sin kx\,dx=\frac{\pi^2(-1)^{k+1}}k+\frac{2((-1)^k-1)}{k^3}$, así que
\[
b_k=\frac{\pi(-1)^{k+1}}{k}+\frac{2\big((-1)^k-1\big)}{\pi k^3}=\begin{cases}\dfrac\pi k-\dfrac4{\pi k^3}&k\text{ impar}\\[6pt]-\dfrac\pi k&k\text{ par.}\end{cases}
\]
En $x=\pi$ hay un salto (por la izquierda vale $\pi^2$ y por la derecha $0$): la serie converge a $\frac{\pi^2}2$.
(Con la convención antigua $\frac{a_0}{\sqrt2}$ de la pauta, el término constante es el mismo, $\frac{\pi^2}6$.)

\sol{5}
$a_0=\frac1{2\pi}\int_{\pi/2}^\pi x\,dx=\frac{3\pi}{16}$. Por partes:
\[
a_k=\frac1\pi\Big[\frac{x\sin kx}k+\frac{\cos kx}{k^2}\Big]_{\pi/2}^{\pi}=\frac1\pi\Big(\frac{(-1)^k-\cos\frac{k\pi}2}{k^2}-\frac{\pi\sin\frac{k\pi}2}{2k}\Big),
\]
\[
b_k=\frac1\pi\Big[-\frac{x\cos kx}k+\frac{\sin kx}{k^2}\Big]_{\pi/2}^{\pi}=\frac{(-1)^{k+1}}k+\frac{\cos\frac{k\pi}2}{2k}-\frac{\sin\frac{k\pi}2}{\pi k^2}.
\]
Por ejemplo, $a_1=-\frac1\pi-\frac12$, $b_1=1-\frac1\pi$, $a_2=\frac1{2\pi}$, $b_2=-\frac34$.
Convergencia: en $x=\frac\pi2$ hay un salto de $0$ a $\frac\pi2$, así que converge a $\frac\pi4$. En $x=\pi$ hay un salto de $\pi$ a $0$ (pues $g(-\pi^+)=0$), así que converge a $\frac\pi2$.

\sol{6}
$q$ es la onda cuadrada de período $T$, así que $f=\frac12+\frac12\cdot\frac4\pi\sum_{k\text{ impar}}\frac1k\sin\frac{2\pi kt}{T}$. Es decir, $a_0=\frac12$ y $b_k=\frac2{\pi k}$ para $k$ impar.
Armónicos: 1000 Hz con $A_1=\frac2\pi\approx0.637$; 3000 Hz con $0.212$; 5000 Hz con $0.127$.

\sol{7}
(a) $A_3=\sqrt{4+16}=\sqrt{20}$ y la potencia es $\frac12\cdot20=10$.
(b) $3^2+\frac12(2^2+1^2)=11.5$. Las oscilaciones aportan $\frac{2.5}{11.5}\approx21.7\%$.
(c) $1-\frac8{\pi^2}\big(1+\frac19+\frac1{25}\big)\approx0.067$, es decir $6.7\%$.
(d) $\frac1{2\pi}\int x^2=\frac{\pi^2}3=\frac{\pi^2}4+\frac12\cdot\frac{16}{\pi^2}S$, con $S=\sum_{\text{impar}}\frac1{k^4}$. Entonces $S=\frac{\pi^4}{96}$. Los pares aportan $\frac1{16}$ del total, así que $\sum_{k\ge1}\frac1{k^4}=\frac{16}{15}S=\frac{\pi^4}{90}$.

\sol{8}
Energías $E_k=\frac12A_k^2$: $0.18,\ 2,\ 0.02,\ 0.72$, con total $2.92$. Conviene guardar los armónicos 2 y 7 (no el 1 y el 2). Se descarta $0.18+0.02=0.2$, y el error relativo es $\sqrt{0.2/2.92}\approx0.26$.

\sol{9}
(a) $f_s/2=500$ Hz. 450 Hz: sin aliasing. $700\mapsto300$ Hz, $1300\mapsto300$ Hz, $2000\mapsto0$ Hz (parece constante).
(b) $|300-250|=50$ Hz: se confunde con la interferencia de la red eléctrica, que ya está en 50 Hz.
(c) $26-24=2$: parece girar hacia adelante a 2 vueltas/s.

\end{document}
